\documentclass[main]{subfiles}
\usepackage[sols]{workbook}
\numbering{ws1-5}

\begin{document}

\ifSubfilesClassLoaded{%
  \chaptermark{\chapterone}
}{}

\section{3D Density Problems} \label{ws1-5}

\begin{framed}
\textbf{Lesson Goals:}
\begin{itemize}
\item You should be able to apply our general slicing strategy to find the total mass/amount in a 3D region with a given density function.
\item You should be able to decide how to draw, label, orient, and slice a 3D region with a given density function, and apply various geometric techniques to approximate the amount within each slice. 
\end{itemize}
\end{framed}

\begin{enumerate}
 

\item  \label{density-cone-1} % From Robin's Spring 2007 in-class worksheet
  \begin{problem}[3in]
    A cone with height $8$ inches and radius $6$ inches is filled with
    flavored slush. When the cup is held upright with the pointed end
    resting on a table, the density of flavoring syrup in the cup varies
    with height above the table.  Suppose $\rho(x)$ gives the number of
    ounces of syrup per cubic inch, where $x$ is the distance from the
    table top.  Write an integral giving the total amount of syrup in the
    cup.
  \end{problem}

  \begin{solution}
    As with the 2-dimensional density problems we did earlier, the key in
    density problems is is always to slice so that the density within each
    slice is approximately constant.  This allows us to reasonably
    approximate the amount (of syrup, in this case) in each slice by
    multiplying the density within the slice by the volume of the slice.  
    
    Here is a diagram of the cone along with the ``$x$-axis'' for this
    problem.  (Since $x$ represents distance from the table top, $x = 0$
    at the tip of the cone and $x = 8$ at the top.)
    \begin{center}
      \includegraphics[height=1.5in]{ws/images/density-cone-1}
    \end{center}
    Since the density varies with height, we will slice the cone parallel
    to the table at a point $x_k=k\Delta x$.  Then, the $k$-th slice looks approximately like a disk
    of thickness $\Delta x=\frac{8}{n}$, and we will let $r_k$ be the radius of the
    top of this slice:
    \begin{center}
      \includegraphics[height=1.5in]{ws/images/density-cone-1-slice}
    \end{center}
    The volume of the $k$-th slice is approximately $\pi r_k^2 \Delta x$.
    However, since we really want to write everything in terms of $x_k$
    and $\Delta x$, we need to write $r_k$ in terms of $x_k$.  To do this,
    we can use similar triangles:
    \begin{center}
      \includegraphics[width=1.5in]{ws/images/density-cone-1-cross-section}
    \end{center}
    This shows us that $\frac{6}{8} = \frac{r_k}{x_k}$, so $r_k =
    \frac{6}{8} x_k = \frac{3}{4} x_k$.  Therefore,
    \begin{align*}
      \textrm{volume of $k$-th slice} &\approx \pi \left( \frac{3}{4} x_k
      \right)^2 \Delta x \\
      & = \frac{9 \pi}{16} x_k^2 \Delta x \textrm{ cubic inches} \\
      \textrm{density of syrup in $k$-th slice} & \approx \rho(x_k)
      \textrm{ ounces per cubic inch}\\
      \textrm{amount of syrup in $k$-th slice} & \approx \rho(x_k) \cdot 
        \frac{9 \pi}{16} x_k^2 \Delta x \\
      & = \frac{9 \pi}{16} x_k^2 \rho(x_k)
    \Delta x \textrm{ ounces}
    \end{align*}
    Summing over all slices, the amount of syrup in the cup is
    approximately $\displaystyle \sum_{k=1}^n \frac{9 \pi}{16} x_k^2
    \rho(x_k) \Delta x$.  Taking the limit as $n \to \infty$, the amount
    of syrup in the cup is exactly \fbox{$\displaystyle \int_0^8 \frac{9
        \pi}{16} x^2 \rho(x)~dx$}.
  \end{solution}

\item  \label{density-concentric-spherical-shells-1}
  % From Robin's Spring 2007 in-class worksheet

  \note{The most important thing is to make sure students know how to
    slice and can describe what a slice looks like.  If students are
    having trouble, remind them that the key is to slice so that the
    density within a slice is approximately constant.  Most students will
    realize that the volume of a spherical shell can be approximated by
    surface area multiplied by thickness; some of these students won't
    remember the formula for surface area, so remind them.}
  
  \begin{problem}[2.8in]
    Suppose the density of a planet is given by the function
    $\displaystyle \rho(r) = \frac{40000}{1 + 0.0001r^3}$ kilograms per
    cubic kilometer, where $r$ is the distance in kilometers from the
    center of the planet. Find the total mass of the planet if its radius
    is $8000$ km. You do not need to evaluate your integral. 
    
    \textit{Hint: You might use the fact that the surface area of a sphere of radius $r$ is $4\pi r^2$}. 
  \end{problem}

  \begin{solution}
    Since the density varies with distance to the center, we should slice
    into concentric spherical shells.  This amounts to slicing the
    interval $[0, 8000]$ into $n$ pieces of equal width $\Delta r=\frac{8000}{n}$, and we
    label $r_0 = 0, \ldots, r_n = 8000$ as usual.  Each shell will have a
    small thickness $\Delta r$, and a good approximation for the volume of
    such a spherical shell is its surface area multiplied by the thickness
    $\Delta r$.

    The $k$-th slice has outer radius $r_k=k\Delta r$, so
    \begin{align*}
      \textrm{surface area of $k$-th slice} & \approx 4 \pi r_k^2 \\
 i     \textrm{volume of $k$-th slice (spherical shell)} & \approx 4 \pi
      r_k^2 \Delta r \\
      \textrm{density of $k$-th slice} & \approx \rho(r_k) \\
      \textrm{mass of $k$-th slice} & \approx 4 \pi r_k^2 \rho(r_k) \Delta
      r
    \end{align*}
    Summing over all slices, the mass of the planet is approximately
    $\displaystyle \sum_{k=1}^n 4 \pi r_k^2 \rho(r_k) \Delta r$.  Taking
    the limit, the mass of the planet is exactly \fbox{$\displaystyle
      \int_0^{8000} 4 \pi r^2 \rho(r)\ dr = \int_0^{8000} 4 \pi r^2 \frac{40000}{1 + 0.0001r^3}\ dr$}.  If we wanted a numerical
    answer, we could integrate using the substitution $u = 1 + 0.0001r^3$.
  \end{solution}
\pspace{\newpage}

\item  \label{density-candle}

  \note{In this problem, it's important to distinguish between $x$ and the
    radius.  Some students will use $x$ to mean both the distance to the
    wick and the distance to the center of the candle.}

  A cylindrical candle of height $50$ mm and radius $12$ mm is formed by
  repeatedly dipping a wick of radius $1$ mm into hot wax and then
  allowing the new layer of wax to dry.  The density of each new layer is
  slightly different, so the density of the candle varies with the
  distance from the wick.  Suppose $\rho(x)$ gives the density in grams per
  cubic mm of the wax, where $x$ measures the distance from the \textit{wick} (in
  mm). 

  \begin{enumerate}
  \item 
    \begin{problem}[2.6in]
      Write an integral giving the mass of the candle.  (The mass of the
      wick is negligible.)
    \end{problem}

    \begin{solution}
      Here is a (rudimentary) sketch of our candle.
      \begin{center}
        \includegraphics[width=1.5in]{ws/images/density-candle}
      \end{center}
      Notice that, since $x$ measures distance to the wick and the wick has
      radius 1, $x = 0$ is on the edge of the wick and $x = 11$ is at the edge of
      the candle.

      We want to slice so that the density within each slice is
      approximately constant.  Since density varies with distance to the
      wick, we should slice into cylindrical shells (paper towel tubes, just
      like what we used in some of the volumes of revolution problems).
      This amounts to slicing the interval $[0, 11]$ into $n$ pieces of
      equal width $\Delta x=\frac{11}{n}$; we label $x_0 = 0, \ldots, x_n = 11$ as usual.

      As usual, we first focus on a typical $k$-th slice. The $k$-th cylindrical shell has height $50$ mm, radius $r_k$, and thickness $\Delta x$. Just as before, we want to express everything in terms of $x$, which means we must express $r_k$ in terms $x_k=k\Delta x$. The outer radius $r_k$ of the $k$-th tube is equal to $x_k + 1$. (It is not just $x_k$, because $x_k$ only measures the distance between the edge of the $k$-th slice and the edge of the \textit{wick}. So we must add 1 to take into account the radius of the wick when measuring the radius of the $k$-th shell).  So:
      \begin{align*}
        \textrm{volume of $k$-th slice} & \approx 2 \pi (r_k) \cdot 50 
        \cdot \Delta x\\
       & =2 \pi (x_k + 1) \cdot 50 
        \cdot \Delta x \\
        & = 100 \pi (x_k + 1) \Delta x \textrm{ cubic mm} \\
        \textrm{density of wax in $k$-th slice} & \approx \rho(x_k) \textrm{
          grams per cubic mm} \\
        \textrm{mass of $k$-th slice} & \approx 100 \pi (x_k + 1)
        \rho(x_k) \Delta x \textrm{ grams}
      \end{align*}
      Summing over all slices, the mass of the candle is approximately
      $\displaystyle \sum_{k=1}^n 100 \pi (x_k + 1) \rho(x_k) \Delta x$.
      Taking the limit as $n \to \infty$, the mass of the candle is exactly 
      \fbox{$\displaystyle \int_0^{11} 100 \pi (x + 1) \rho(x)~dx$}.
    \end{solution}

  \item

    \note{If they have the wrong answer to the previous part, they should
      find that out here.}
    
    \begin{problem}[1.5in]
      Show that, if $\rho(x) = \frac{1}{x + 1}$, then the mass of the
      candle is $1100 \pi$ grams.
    \end{problem}

    \begin{solution}
      We already showed that the mass of the candle is $\displaystyle
      \int_0^{11} 100 \pi (x + 1) \rho(x)~dx$.  If $\rho(x) = \frac{1}{x +
        1}$, then this is equal to $\displaystyle \int_0^{11} 100 \pi~dx =
      1100 \pi$.
    \end{solution}
  \end{enumerate}

  \pspace{\newpage}
  \item  Lake Baikal (also known as the Blue eye of Siberia) is the deepest lake
  in the world; its depth is approximately $1500$ meters. Baikalsk Pulp
  and Paper Mill was constructed directly on the shore line. It bleaches
  its paper with chlorine and discharges the waste into Baikal. Despite
  numerous protests, the Mill is still in production.  The waste
  accumulates at the bottom of the lake. The density of the waste in the
  lake varies with depth and is given by function $\rho(\ell)$ kg/m$^3$,
  where $\ell$ represents depth (measured in meters).
  
  Let us suppose that the shape of the lake can be described as the solid
  of revolution obtained by rotating the region in the first quadrant
  bounded by $y = 1500$, $y = x^2$, and the $y$-axis about the $y$-axis.
  
  \begin{enumerate}
  \item
    \begin{problem}
      Write a general Riemann sum that approximates the total amount of
      waste in the lake. 
    \end{problem}
\pspace{\vfill}
    \begin{solution}
      Here is a diagram:
      \begin{center}
        \includegraphics[height=1.5in]{ws/images/density-baikal}
      \end{center}
      Remember that the depth $\ell$ is 0 at the surface of the lake and
      1500 at the bottom of the lake.  Therefore, the relationship between
      $\ell$ and $y$ is $\ell = 1500 - y$.
      
      Since we want to make slices in which the density of waste is
      approximately constant, we need to slice parallel to the surface of
      the lake.  We can either think of slicing $\ell$ or think of slicing
      $y$; in this solution, we'll slice $\ell$.  We slice the interval
      $[0, 1500]$ into $n$ pieces of equal height $\Delta \ell$ and label
      the endpoints $\ell_0 = 0, \ldots, \ell_n = 1500$ as usual.  Let's
      focus on the $k$-th slice and label the top right point in the slice
      $(x_k, y_k)$:
      \begin{center}
        \includegraphics[height=1.5in]{ws/images/density-baikal-slice}
      \end{center}      
      This slice can be approximated by a disk of radius $x_k$ and
      thickness $\Delta \ell$; therefore, its volume is approximately $\pi
      x_k^2 \Delta \ell$, and the amount of waste in this slice is
      approximately $\pi x_k^2 \rho(\ell_k) \Delta \ell$.

      In order to have a Riemann sum that makes sense, we need to write
      $x_k$ in terms of $\ell_k$.  We know that $(x_k, y_k)$ is on the
      curve $y = x^2$, so $y_k = x_k^2$.  That is, $x_k = \sqrt{y_k}$.  We
      also know that the relationship between $\ell_k$ and $y_k$ is
      $\ell_k = 1500 - y_k$, so $y_k = 1500 - \ell_k$ and $x_k =
      \sqrt{1500 - \ell_k}$.

      So, the amount of waste in the $k$-th slice is approximately $\pi
      (1500 - \ell_k) \rho(\ell_k) \Delta \ell$.  Summing over all $n$
      slices gives us the Riemann sum approximation $\boxed{\sum_{k=1}^n
        \pi (1500 - \ell_k) \rho(\ell_k) \Delta \ell}$.
    \end{solution}

  \item
    \begin{problem}
      By taking the appropriate limit of the Riemann sum in (a), arrive at
      an integral that gives the exact amount of waste in the lake.
    \end{problem}
\pspace{\vfill}
    \begin{solution}
      The integral is $\displaystyle \int_0^{1500} \pi (1500 - \ell)
      \rho(\ell)~d\ell$.
    \end{solution}
    
  \end{enumerate}
\pspace{\newpage}
  \item  \label{density-muffin}
  
  \begin{problem}
    We can model a muffin as a solid of revolution, obtained by rotating
    the following region about the $y$-axis.  Due to a poor recipe, the
    chocolate chips in our muffin tend to sink to the bottom.  The amount
    of chocolate in the muffin is given by $\rho(y) = 5 - y$ grams per
    cubic inch, where $y$ represents the distance to the bottom of the
    muffin.  Write an integral that gives the total amount of chocolate in the muffin.

    \includegraphics[width=1.75in]{ws/images/density-muffin-1}
  \end{problem}
\pspace{\vfill}
  \begin{solution}
    Note that, in the picture, $y$ does indeed represent distance to the
    bottom of the muffin.  Since the chocolate density varies with the
    distance to the bottom of the muffin, we must slice parallel to the
    bottom of the muffin.  Each slice is approximately a disk.  Here are
    two representative slices:
    \begin{center}
      \includegraphics[width=2in]{ws/images/density-muffin-2}
    \end{center}
    Notice that these two slices have different descriptions: the top one
    should be described using the curve $y = 3 - \frac{x^2}{4}$, while the
    bottom one should be described using the curve $y = 2x - 2$.  So, we
    should really consider the top and bottom parts of the muffin
    separately.

    Let's first focus on the bottom part of the muffin, which we get by
    rotating the following region about the $y$-axis: 
    \begin{center}
      \includegraphics[width=2in]{ws/images/density-muffin-bottom}
    \end{center}
    As we said, we will use horizontal slices, so we slice the interval
    $[0, 2]$ (on the $y$-axis) into $n$ pieces of equal height $\Delta
    y$.  We label $y_0 = 0, \ldots, y_n = 2$ as usual and focus on the
    $k$-th slice:
    \begin{center}
      \includegraphics[width=2in]{ws/images/density-muffin-bottom-slice}
    \end{center}
    The equation of the line $y = 2x - 2$ can be rewritten as $x = \frac{y
      + 2}{2}$.  If we rotate the $k$-th slice of this region, we get a
    slice that is approximately a disk of radius $\frac{y_k + 2}{2}$ and
    thickness $\Delta y$, so the volume of the $k$-th slice of the muffin
    bottom is approximately $\pi \left( \frac{y_k + 2}{2} \right)^2 \Delta
    y$.  The density of chocolate in this slice is approximately
    $\rho(y_k)$, so the amount of chocolate in this slice is approximately
    $\rho(y_k) \cdot \pi \left( \frac{y_k + 2}{2} \right)^2 \Delta y = \pi
    \left( \frac{y_k + 2}{2} \right)^2 \rho(y_k) \Delta y$.  Adding these
    up and taking the limit gives us the integral $\displaystyle \int_0^2
    \pi \left( \frac{y + 2}{2} \right)^2 \rho(y)~dy$ for the amount of
    chocolate in the bottom part of the muffin.

    To deal with the top part of the muffin, we use exactly the same
    reasoning.  We get the top part of the muffin by rotating this region:
    \begin{center}
      \includegraphics[width=2in]{ws/images/density-muffin-top}
    \end{center}
    We need to solve $y = 3 - \frac{x^2}{4}$ for $x$:
    \begin{align*}
      y & = 3 - \frac{x^2}{4} \\
      \frac{x^2}{4} & = 3 - y \\
      x^2 & = 4 (3 - y) \\
      x & = 2 \sqrt{3 - y}
    \end{align*}
    So, the $k$-th slice is approximately a disk of radius $2 \sqrt{3 -
      y_k}$ and thickness $\Delta y$, which means its volume is
    approximately $\pi (2 \sqrt{3 - y_k})^2 \Delta y = 4 \pi (3 - y_k)
    \Delta y$.  Therefore, the amount of chocolate in this slice is
    approximately $4 \pi (3 - y_k) \rho(y_k) \Delta y$.  Adding these up
    and taking the limit gives us the integral $\displaystyle \int_2^3 4
    \pi (3 - y) \rho(y)~dy$ for the amount of chocolate in the top part
    of the muffin.

    So, the total amount of chocolate in the muffin is
    \fbox{$\displaystyle \int_0^2 \pi \left( \frac{y + 2}{2} \right)^2
      \rho(y)~dy + \int_2^3 4 \pi (3 - y) \rho(y)~dy$.}
  \end{solution}

  \item  \begin{problem}[4in]
    A cylindrical sewage pipe $30$ m long and $2$ m in diameter is lying
    on the ground. The density of the refuse in the pipe is given by
    $\rho(x)$ g/m$^3$ where $x$ is the distance from the ground in
    meters. The refuse in the pipe reaches a height of $1.5$ m.  Write an
    integral that gives the mass of refuse in the pipe.
  \end{problem}
\pspace{\vfill}
  \begin{solution}
    Here is a picture of a cross-section of the pipe.  Since $x$ measures
    distance from the ground in meters, we've made it the vertical axis,
    and we've given the horizontal axis a different name ($u$): 
    \begin{center}
      \includegraphics[width=1.5in]{ws/images/sewage-pipe}
    \end{center}    
    We want to make slices in which the density of refuse is approximately
    constant, so we need to slice parallel to the ground.  The bottom of
    the pipe is where $x = 0$, and the top of the refuse is where $x =
    1.5$, so we slice the interval $[0, 1.5]$ into $n$ equal pieces of
    height $\Delta x$, label $x_0 = 0, \ldots, x_n = 1.5$ as usual, and
    focus on the $k$-th slice: 
    \begin{center}
      \includegraphics[width=1.5in]{ws/images/sewage-pipe-slice}
    \end{center}        
    The $k$-th slice is approximately a rectangular slab of length 30 m
    and thickness $\Delta x$.  Its width is the horizontal distance from
    the left side of the circle to the right side of the circle.  The
    easiest way to calculate this width is using the equation of this
    circle, which is $(x - 1)^2 + u^2 = 1$.  Solving for $u$, we get $u =
    \pm \sqrt{1 - (x - 1)^2}$.  The right side is $u = \sqrt{1 - (x -
      1)^2}$ and the left side is $u = -\sqrt{1 - (x - 1)^2}$, so the
    horizontal distance is $\sqrt{1 - (x_k - 1)^2} - [-\sqrt{1 - (x_k -
      1)^2}] = 2 \sqrt{1 - (x_k - 1)^2}$.  Therefore, the volume of this
    slab is approximately $30 \cdot 2 \sqrt{1 - (x_k - 1)^2} \cdot \Delta
    x = 60 \sqrt{1 - (x_k - 1)^2} \Delta x$.

    The mass of refuse in this slice is approximately volume times
    density, or $60 \sqrt{1 - (x_k - 1)^2} \rho(x_k) \Delta x$.  Summing
    over all slices and taking the limit as the number of slices tends to
    $\infty$, we get the integral \fbox{$\displaystyle \int_0^{1.5} 60
      \sqrt{1 - (x - 1)^2} \rho(x)~dx$}.
  \end{solution}


\end{enumerate}




\end{document}